\chapter{Droites, projections et cercles}
\section*{Exercice 18.1 -- Droite normale}
La condition est $3(x-2)+2(y+1)=0$, soit $3x+2y-4=0$. Son vecteur directeur peut être $(-2,3)$ ; le produit avec $(3,2)$ vaut $0$.

\section*{Exercice 18.2 -- Projeté}
Le coefficient de projection vaut $(9+16-5)/(9+16)=4/5$. Donc $H=(3,4)-(4/5)(3,4)=(3/5,4/5)$ et $MH=20/5=4$. Vérifier $3(3/5)+4(4/5)-5=0$ contrôle l'appartenance à la droite.

\section*{Exercice 18.3 -- Reconnaître un cercle}
Compléter les carrés donne $(x+1)^2+(y-2)^2=5-c$. Si $c<5$, c'est un cercle de centre $(-1,2)$ et rayon $\sqrt{5-c}$. Si $c=5$, un point ; si $c>5$, l'ensemble vide.

\section*{Exercice 18.4 -- Intersections}
Substituer donne $2x^2=10$, donc $x=\pm\sqrt5$. Comme $y=x$, les points sont $(\sqrt5,\sqrt5)$ et $(-\sqrt5,-\sqrt5)$. On vérifie les deux équations pour chacun.

\section*{Exercice 18.5 -- Tangente}
Le rayon $(3,4)$ est normal, donc l'équation est $3(x-3)+4(y-4)=0$, soit $3x+4y=25$. Pour un point $P$ sur cette droite, Pythagore donne $OP^2=25+TP^2$. L'égalité $OP^2=25$ impose $P=T$.

\section*{Exercice 18.6 -- Foyer et directrice}
L'égalité des distances est $\sqrt{x^2+(y-2)^2}=|y+2|$. En élevant au carré, on obtient $x^2=8y$. Avec $y=2$, $x^2=16$, donc les points $(-4,2)$ et $(4,2)$. Les distances étant positives, l'élévation au carré est réversible.
